\documentclass[12pt]{amsart}
\usepackage[letterpaper,margin=1in]{geometry}
%%
%%%%%%%%%%%%%%%%%%%%%%%%%%%%%%%%%%%%%%%%%%%%%%%%%%%%%%%%%%%%%%%%%%%%%%%%%%%%
\usepackage{listings}
% \lstdefinestyle{pystyle}{
%        language=Python,
%        basicstyle=\ttfamily,
%        keywordstyle=\color{blue},
%        stringstyle=\color{red},
%        commentstyle=\color{green},
%        numbers=left,
%        numberstyle=\tiny\color{gray},
%        stepnumber=1,
%        showstringspaces=false
%    }
%%%%%%%%%%%%%%%%%%%%%%%%%%%%%%%%%%%%%%%%%%%%%%%%%%%%%%%%%%%%%%%%%%%%%%%%%%%%

\usepackage{amsmath,amssymb,amsfonts}
\usepackage{bm}
\usepackage{bbm}
\usepackage[config,labelfont={rm,rm},textfont=rm]{caption,subfig}

\usepackage{amsmath,amssymb,amsthm,latexsym,mathrsfs}

\usepackage{color,graphicx,colortbl}

\usepackage[latin1]{inputenc}
\usepackage[english]{babel}

\theoremstyle{definition}
\newtheorem{theorem}{Theorem}[section]
\newtheorem{lemma}[theorem]{Lemma}
\newtheorem{problem}[theorem]{Problem}
\newtheorem{proposition}[theorem]{Proposition}
\newtheorem{corollary}[theorem]{Corollary}
\newtheorem{definition}[theorem]{Definition}
\newtheorem{algorithm}[theorem]{Algorithm}
\newtheorem{example}[theorem]{Example}
\newtheorem{remark}[theorem]{Remark}
\newtheorem{exercise}[theorem]{Exercise}

\renewcommand{\labelenumi}{\textbf{\arabic{enumi}.}}
\renewcommand{\labelenumii}{\textbf{\arabic{enumi}\alph{enumii}.}}
\everymath{\displaystyle}

\title{An idea for solution. Question(s) for Zhilan \& John}
\author{\ldots \and Yuliya \and Ludmil}

\begin{document}

\maketitle

\begin{abstract}
On the minimizng spectral radius. 
\end{abstract}

\section{A few easy question(s)}
\begin{itemize}
\item Let us consider the case $n=3$, $m=2$ and  $\bm F$ defined as follows (this is from Zhilan's and John's  writeup):
\begin{equation}
\bm F = \begin{pmatrix}
F_{1111} & F_{1112} & F_{1113} &  F_{1121} &  F_{1122} & F_{1123} \\
F_{1211} & F_{1212} & F_{1213} & F_{1221} &  F_{1222} & F_{1223} \\
F_{1311} & F_{1312} & F_{1313} & F_{1321} &  F_{1322} & F_{1323} \\
F_{2111} & F_{2112} & F_{2113} & F_{2121} &  F_{2122} & F_{2123} \\
F_{2211} & F_{2212} & F_{2213} & F_{2221} &  F_{2222} & F_{2223} \\
F_{2311} & F_{2312} & F_{2313} & F_{2321} &  F_{2322} & F_{2323}
\end{pmatrix}
\end{equation}
\item One question is about the definition of $\bm F_{ijpq}$ (which in the manuscript is said to be $8\times 8$. {\color{red} It represents the interactions between susceptible people of age $j$ at location $i$ with infected people of age $q$ at location $p$}). 
Let us fix $(i,j,p,q)$ and look at the definition of $\bm F_{ijpq}$:
\begin{equation}\label{f-ijpq-1}
\bm{F}_{ijpq}=
\begin{pmatrix}
0 & 0 &  Y_{ij}\Psi_{qI}\eta_{a}c_{ijpq} & \ldots & Y_{ij}\Psi_{qI}\eta_{m}c_{ijpq}\\
0 & 0 &  Z_{ij}\Psi_{qI}\eta_{a}c_{ijpq} & \ldots & Z_{ij}\Psi_{qI}\eta_{m}c_{ijpq}\\
0 & 0 &  0 & \ldots & 0\\
\vdots & \vdots &  \vdots & \ldots & \vdots\\
0 & 0 &  0 & \ldots & 0
\end{pmatrix}
\end{equation}
\end{itemize}

We change the indices here to $K=(j-1)*n+i$, $i=1:n$, $j=1:m$ \textcolor{blue}{YG: Here is $n$ and $m$ are yet again switched...}  and $L=(q-1)*n+p$, $p=1:n$, $q=1:m$ then we have a "two-dimensional" matrix $F_{KL}$.

For the case of $m=4$ and $n=4$,
the first row  of block matrices of  $F$  is 
\tiny{
\[
(F_{1111} \ \ F_{1112} \ \ F_{1113} \  \ F_{1114} \  \
F_{1121} \ \  F_{1122} \ \ F_{1123}  \ \ F_{1124}   \ \
F_{1131} \ \   F_{1132} \ \  F_{1133} \ \  F_{1134} \ \
F_{1141} \ \   F_{1142} \ \  F_{1143} \ \  F_{1144}).
\]
}
{\color{red}(ZF note: it is a row vector of block matrices with dimension 1 by 16. It describes the interactions between people of age $j=1$ at location $i=1$ with infected people of age $q=1,2,3,4$ at location $p=1,2,3,4$.)}

The first row  of   $F_{ijpq}$  is
(consider the simpler case when $\eta$ and $\Psi$ are independent of age or location)
\tiny{
\begin{equation}\label{F-def-1}
\begin{pmatrix} 
0 & 0 & Y_{ij}\Psi_{qI}\eta_{a}c_{ijpq} & Y_{ij}\Psi_{qI}\eta_{s}c_{ijpq} &Y_{ij}\Psi_{qI}c_{ijpq} 
&Y_{ij}\Psi_{qI}\eta_{h}c_{ijpq}  &
Y_{ij}\Psi_{qI}\eta_{pm}c_{ijpq}  & Y_{ij}\Psi_{qI}\eta_{m}c_{ijpq}  
\end{pmatrix}.
\end{equation}
} %end tiny
So, row one of $F_{1111}$ is
\tiny{
\[
\begin{pmatrix}
0 & 0 & Y_{11}\Psi_{qI}\eta_{a}c_{1111} &
Y_{11}\Psi_{qI}\eta_{s}c_{1111} &
Y_{11}\Psi_{qI}c_{1111} &
Y_{11}\Psi_{qI}\eta_{h}c_{1111} &
Y_{11}\Psi_{qI}\eta_{pm}c_{1111} &
Y_{11}\Psi_{qI}\eta_{m}c_{1111}
\end{pmatrix},
\]
}
Row one of $F_{1112}$ is
\tiny{\[
\begin{pmatrix}
0 & 0 & Y_{11}\Psi_{qI}\eta_{a}c_{1112} &
Y_{11}\Psi_{qI}\eta_{s}c_{1112} &
Y_{11}\Psi_{qI}c_{1112} &
Y_{11}\Psi_{qI}\eta_{h}c_{1112} &
Y_{11}\Psi_{qI}\eta_{pm}c_{1112} &
Y_{11}\Psi_{qI}\eta_{m}c_{1112}
\end{pmatrix},
\]
}

\vdots
\noindent
Row one of $F_{1144}$ is
\tiny{\[
\begin{pmatrix}
0 & 0 & Y_{11}\Psi_{qI}\eta_{a}c_{1144} &
Y_{11}\Psi_{qI}\eta_{s}c_{1144} &
Y_{11}\Psi_{qI}c_{1144} &
Y_{11}\Psi_{qI}\eta_{h}c_{1144} &
Y_{11}\Psi_{qI}\eta_{pm}c_{1144} &
Y_{11}\Psi_{qI}\eta_{m}c_{1144}
\end{pmatrix}.
\]
}

{\color{red}(ZF note: It might be helpful to illustrate here the structure of the $F$ matrix. I have no idea how to type this $F$ matrix using matrix command in LaTeX because the page width does not allow 16 elements in a single row.
The Mathematica command for this F matrix is: 

mF = Table[
   Table[0, \{r, 1, 8\}, \{s, 1, 8\}], \{i, 1, m\}, \{p, 1, m\}, \{j, 1, 
    n\}, \{q, 1, n\}]
    
where Table[0, \{r, 1, 8\}, \{s, 1, 8\}] generates a matrix of 8 by 8 for each $F_{ijpq}$. The matrix structure  from this command is as shown in Figure \ref{fig:StructureF}. There are also 4 by 4 block of $F_{ijpq}$ matrices separating the 4 locations.
)}

The full matrix of $F$:
\tiny{
 \[
Row 1: (F_{1111} \ \ F_{1112} \ \ F_{1113} \  \ F_{1114} \  \
F_{1121} \ \  F_{1122} \ \ F_{1123}  \ \ F_{1124}   \ \
F_{1131} \ \   F_{1132} \ \  F_{1133} \ \  F_{1134} \ \
F_{1141} \ \   F_{1142} \ \  F_{1143} \ \  F_{1144})
\]
\[
Row 2: (F_{1211} \ \ F_{1212} \ \ F_{1213} \  \ F_{1214} \  \
F_{1221} \ \  F_{1222} \ \ F_{1223}  \ \ F_{1224}   \ \
F_{1231} \ \   F_{1232} \ \  F_{1233} \ \  F_{1234} \ \
F_{1241} \ \   F_{1242} \ \  F_{1243} \ \  F_{1244})
\]
\[
Row 3: (F_{1311} \ \ F_{1312} \ \ F_{1313} \  \ F_{1314} \  \
F_{1321} \ \  F_{1322} \ \ F_{1323}  \ \ F_{1324}   \ \
F_{1331} \ \   F_{1332} \ \  F_{1333} \ \  F_{1334} \ \
F_{1341} \ \   F_{1342} \ \  F_{1343} \ \  F_{1344})
\]

\quad \vdots
\[
Row 16: (F_{4411} \ \ F_{4412} \ \ F_{4413} \  \ F_{4414} \  \
F_{4421} \ \  F_{4422} \ \ F_{4423}  \ \ F_{4424}   \ \
F_{4431} \ \   F_{4432} \ \  F_{4433} \ \  F_{4434} \ \
F_{4441} \ \   F_{4442} \ \  F_{4443} \ \  F_{4444})
\]
}
{\color{magenta}
LZ: No problem with the form of $F$. If the definition of $F_{ijpq}$ in \eqref{f-ijpq-1} is correct, then this is all thats needed. The matrix $c_{ijpq}$ has 16 elements per row and $\eta$ in~\eqref{f-ijpq-1} has ($6$ plus $2$ zeros) so $8$ elements. In total then should be $128$ elements per row in $F$ (if we join all blocks). Many are zero but we don't care abbout this at the moment. 

Another note:Anyway, we are trying to vectorize the formula for $F$, that is, to write it as a tensor product of several matrices. This is possible, we believe we can have some formula by new years coming to the west coast tomorrow\footnote{today is: \today} :). 

ZF: Yes, F is 128 by 128. It is not easy to type those matrices in LaTeX form due to the large dimension. Mathematica allows easier ways to manipulate the matrix operations.

Same is true for \verb|Python| (LZ) Can Mathematica output a file with the matrix (not in latex format, just like the $C$ matrix in a comma separated text format or even $128\times 128$ excel table.

ZF: We can have the Mathematica output in the form as shown in Figue \ref{fig:StructureF}.
)
} %end magenta
\section{A possible path to solution. Further easy questions\label{s:path}}
\renewcommand{\labelenumi}{\textbf{\arabic{enumi}.}}
\renewcommand{\labelenumii}{\textbf{\alph{enumii}.}}

\subsection{Notation} 
In the following $\bm G\ge 0$, $\bm g\ge 0$ means that the elements of a matrix $\bm G$ or a vector $\bm g$ are non-negative. 

\subsection{Possible path to the solution} We recall here the Perron-Frobenius theorem about non-negative and irreducible matrices because this is an essential result to be used. The proof of this famous theorem is found at many places, for example, check Varga~\cite[Theorem 2.7]{2000Varga-aa}. Please, pay particular attention to item 3. in the statement of the theorem. 
\begin{theorem}[Perron-Frobenius]
  Let $\bm A\geq 0$ be a non-negative and irreducible matrix. Then
  \begin{enumerate}
  \item $\bm A$ has positive real eigenvalue equal to $\rho(\bm A)$.
  \item The corresponding eigenvector has nonnegative elements.
  \item $\rho(\bm A)$ increases when any entry of $\bm A$ increases.
  \item $\rho(\bm A)$ is a simple (not repeating) eigenvalue of $\bm A$.
  \end{enumerate}
\end{theorem}

We now state a helpful result with an obvious proof. 
\begin{lemma}\label{l-1}
Let $\bm A\ge 0$, $\bm B\ge 0$, and $\bm C\ge 0$ be nonnegative matrices. If $\bm B$ and $\bm C$  satisfy the inequality $\bm B \ge \bm C$ (to be understood element-wise),  then 
\begin{equation}
\bm{A}\bm B\ge \bm{A}\bm C,\qquad\mbox{where the inequality holds element-wise.}
\end{equation}
\end{lemma} 
\begin{proof} Since $(\bm B - \bm C)\ge 0$ and $\bm A\ge 0$, it is obvious that $\bm A(\bm B - \bm C)\ge 0$ which is what we wanted to show. 
\end{proof}

Here is a possible path to the solution of the problem in hand:
\begin{enumerate}
\item As discussed at NSF: If $(z,\bm\varphi)$ is eigenpair for $\bm F\bm V^{-1}$, then $(z,\bm \psi)$ is 
an eigenpair for $\bm V^{-1}\bm F$ with $\bm \psi=\bm V\bm \varphi$. \textcolor{blue}{YG: should it be $\bm \psi=\bm V^{-1}\bm \varphi$??} So the eigenvalues are the same and we will focus on minimizing 
the dominant eigenvalue $z=\rho\left(\bm V^{-1} \bm F\right)$.  
\item Let $\bm H(\bm \chi)=\bm{V}^{-1}\bm F(\bm \chi)$ be the matrix for which we would like to minimize the dominant eigenvalue with respect to $\bm \chi$. Since $\bm V^{-1}\ge 0$ (this is easy to prove), from Lemma~\ref{l-1} and the Perron-Frobenius Theorem (item 3.) we have the following result.
\begin{theorem}
Let  $\widehat{\bm \chi}$ and $\widetilde{\bm \chi}$ be two values of the parameters satisfying the constraints. We then have
\begin{equation}\label{monotonicity} 
\mbox{If} \quad 
\bm F(\widehat{\bm\chi})\ge \bm F(\widetilde{\bm\chi})\quad\mbox{then}\quad
\rho\left(\bm H(\widehat{\bm\chi})\right) \ge \rho\left(\bm H(\widetilde{\bm\chi})\right) 
\end{equation}
\end{theorem}
\begin{proof}
The proof follows from Lemma~\ref{l-1} and the Perron-Frobenius Theorem.
\end{proof}
\item Basically, we need to minimize the entries of $\bm F$ with respect to $\bm{\chi}$ satisfying the constraints. That is it.
\item From this observation, it is clear that that {\bf we do not need to know the dominant eigenvalue or the gradient with respect to $\bm \chi$ to find the optimal $\bm \chi$}.
\end{enumerate}


\section{An equivalent minimization problem}

We first write an equivalent form of $\bm F_{ijpq}$  which may be useful later when we minimize the entries of $\bm F(\bm \chi)$.
Since $(i,j,p,q)$ are fixed for the moment, there are things we can take out in front of the matrix and also make the notation a bit more compact. We have,
\begin{equation}\label{f-ijpq-2}
\bm{F}_{ijpq}=
\Psi_{pqI}c_{ijpq}
\begin{pmatrix}
Y_{ij} & 0 &  0 & \ldots & 0\\
Z_{ij} & 0  & 0 & \ldots & 0\\
0 & 0 &  1 & \ldots & 0\\
\vdots & \vdots &  \vdots & \ldots & \vdots\\
0 & 0 &  0 & \ldots & 1
\end{pmatrix}
\begin{pmatrix}
0 & 0 &  \eta_{a} & \ldots & \eta_{m}\\
0 & 0 &  0 & \ldots & 0\\
0 & 0 &  0 & \ldots & 0\\
\vdots & \vdots &  \vdots & \ldots & \vdots\\
0 & 0 &  0 & \ldots & 0
\end{pmatrix}
\end{equation}

As is discussed in \S\ref{sec:answers} the {\bf only} entries of $F_{ijpq}$ which depend on $\bm \chi$ are $Y_{ij}$ and $Z_{ij}$, $i=1:m$, $j=1:n$. 

Next, we can apply the same argument as in Lemma~\ref{l-1} to conclude that to minimize all entries of $F_{ijpq}$ we need to only minimize 
$Y_{ij}(\bm \chi)$ and $Z_{ij}(\bm \chi)$. 

\noindent\hrulefill

\centerline{{\bf QUESTIONS}}

\noindent\hrulefill

{\bf What are the actual definitions of $Y_{ij}$ and $Z_{ij}$ {\color{red}(ZF: Note that a location-dependent parameter $s_i$ has been included. Also, the parameter $a$ is now dependent of both age $j$ and location $i$)}. 
{\color{magenta} (YGLZ: We got this; we were wondering how a suddenly became a matrix; it is now clear.)}
}% {Are they as follows:}

\begin{equation}\label{yz-ij-def}
\begin{aligned}
Y_{ij} &= \Psi_{jS}s_ia_{ij}\beta_j \frac{S_{ij}^*}{N_{ij}}, \quad i=1:m,\; j=1:n\\
Z_{ij} &= \Psi_{jS} s_i a_{ij}\beta_j \frac{(1-\varepsilon_{1})V_{ij}^*+(1-\varepsilon_{3})P_{ij}^*}{N_{ij}}.
\end{aligned}
\end{equation}
Here, we have denoted
\begin{equation}\label{si-def}
\begin{aligned}
\frac{S_{ij}^*}{N_{ij}}&= \left\lbrack
1+\frac{\chi_{1,ij}}{\omega_v+\varepsilon_2\chi_{2,ij}}
\left(1+\frac{\varepsilon_{2}\chi_{2,ij}}{\omega}+
\frac{\varepsilon_{2}\chi_{2,ij}}{\omega}\left\lbrack
1+\frac{\varepsilon_{3}\chi_{3,ij}}{\omega}
\right\rbrack
\right)
\right\rbrack^{-1}\\
\frac{V_{ij}^*}{N_{ij}}&=  \frac{\chi_{1,ij}}{\omega_{v}+\varepsilon_2\chi_{2,ij}} \frac{S_{ij}^*}{N_{ij}}\\
\frac{P_{ij}^*}{N_{ij}}&= \frac{\varepsilon_{2}\chi_{2,ij}}{\omega}\frac{V_{ij}^*}{N_{ij}}=
\frac{\varepsilon_{2}\chi_{2,ij}}{\omega}
\frac{\chi_{1,ij}}{\omega_{v}+\varepsilon_2\chi_{2,ij}} \frac{S_{ij}^*}{N_{ij}} \\
 \frac{R_{ij}^*}{N^*_{ij}}&=
\frac{\chi_{1ij}}{\omega_v+\varepsilon_2 \chi_{2ij}}
\frac{\varepsilon_2 \chi_{2ij}}{\omega}
\Big(1+\frac{\varepsilon_3 \chi_{3ij}}{\omega}\Big)
\frac{S_{ij}^*}{N^*_{ij}}.
\end{aligned}
\end{equation}

We also have $3m$ constraints. Let $S_{ij}(T-1)$, $V_{ij}(T-1)$, and 
$P_{ij}(T-1)$ denote the numbers of individuals in the susceptible $(S)$, vaccinated ($V$), and $(P)$ classes {\color{red} at the end of week $T-1$} (which are the solutions of the ODEs by using the $\chi$ values in the week $T-1$). We need to find optimal $\chi$ values for week $T$. Then the constraints are:
\begin{equation}\label{eq:c}
\begin{array}{l}
\sum_{j=1}^n \chi_{1ij} S_{ij}(T-1) = b_{1i}, \quad i=1:m, \\
\sum_{j=1}^n \chi_{2ij} V_{ij}(T-1) = b_{2i}, \quad i=1:m, \\
\sum_{j=1}^n \chi_{3ij} P_{ij}(T-1) = b_{3i}, \quad i=1:m, 
\end{array}
\end{equation}
where $b_{ki}$ $(k=1, 2, 3)$ are the total available vaccine doses of type $k$ at location $i$ for week $T$. 

{\bf What is the relation between 
$(\chi_{1,ij},\chi_{2,ij},\chi_{3,ij})$ (in the definitions and basically in equations~\eqref{yz-ij-def}--\eqref{eq:c}) and $\chi_{ij}$ (in the constraint)? }

{\bf Answer:} The subscript $k$ represents the vaccine types (e.g., $k=1$ is the primary dose, and $k=2$ and 3 are booster doses.

\section{Parameter values}
For the parameter values of $(a_{ij})$,
$(\beta_{ij})$, $(N_{ij})$,
$(\chi_{1ij})$, $(\chi_{2ij})$, $(\chi_{3ij})$,
and the $c=(c_{ijpq})$ matrix, see the data sets.

$\epsilon_1=\epsilon_2=
\epsilon_3=0.8$;

%$(N_{ij})=
%\begin{pmatrix}
%12861646 & 14862307 & 14294150 & 13831766 %\\
%17039947& 18046593& 
%16847952& 16382252 \\
%31775045 & 34108261 & 
%31760261 & 29018710\\
%19577717 & 22475544& 
%19374289 & 17227215
%\end{pmatrix};
%$ 


$\omega=\omega_v=x \frac{1}{182.5}$, where $x$ is the time unit in days.
The value for $x=7$ (days). John has the number of available vaccine doses for each week, so the optimization is done weekly. After each week, the number of susceptible individuals $(S_{ij})$ and number of vaccinated people $(V_{ij})$  will change, so they need to be updated for the optimal solution in the following week.

$$
\vec{\eta} = (\eta_{a}, \eta_{ps}, \eta_{s}, \eta_{h}, \eta_{pm}, \eta_{m}) = 
(0.5, 1.25, 1, 0.1, 1, 0.75); 
$$ ({\color{magenta}\bf Could be wrong order};  independent of age or location.)

This below is what is in the Mathematica NB:
$$
\vec{\eta} = (\eta_{a}, \eta_{ps}, \eta_{pm}, \eta_{s}, \eta_{m}, \eta_{h}) = 
(0.5, 1.25, 1, 1, 0.75,0.1); 
$$ (independent of age or location.)


$
\gamma_a = (1/7) x; \gamma_s = (1/5) x; \gamma_m = (1/5) x; \gamma_h = 
 (1/10) x; 
k=(1/5) x; \xi=(1/2) x;
$

$(p_j)=(0.971, 0.894, 0.655, 0.289); 
(q_j)=(0.985, 0.947, 0.827, 0.644); 
%\gamma_{s \to h}=
%\gamma_{m\to h}=(1/2) x; 
%\delta_s=0.5; \delta_m=0.2;
$

$(\Psi_{jS})=
(0.823, 0.824, 0.829, 0.824);
(\Psi_{qI})=(0.985, 0.947, 0.827, 0.644);$

$(s_j)=(0.7024, 0.7350, 0.7178, 0.6536).$
%The mixing matrix $c=(c_{ijpq})$
%as a 4 by 4 block matrix or
%a 16 by 16 matrix given below:

%\begin{figure}[htb!]
%    \centering
%    \includegraphics[width=400pt]{c_Matrix.pdf}
%    \caption{The mixing matrix $c$ whose %block matrices are $c_{ijpq}$.}
%    \label{fig:diag_RSV}
%\end{figure}
%Or as a 16 by 16 matrix:

%\begin{figure}[htb!]
%    \centering
%    \includegraphics[width=400pt]%{c_matrix_v2.pdf}
%    \caption{The 16 by 16 mixing matrix %$c$.}
%    \label{fig:diag_RSV}
%\end{figure}


\begin{thebibliography}{9}
\bibitem{2000Varga-aa}
Richard~S. Varga.
\newblock {\em Matrix iterative analysis}, volume~27 of {\em Springer Series in Computational Mathematics}.
\newblock Springer-Verlag, Berlin, expanded edition, 2000.

\end{thebibliography}


%%%%%%%%%%%%%%%%%%%%%%%%%%%%%%%%%%%%%

\newpage
\section{Answered\label{sec:answers}}
\begin{itemize}
\item If $\bm V$ is defined as in the document you gave us then we have that $\bm V$ is non-zero for $i=p$, and zero for $i\neq p$ (but nothing is said about $j$ and $q$). The definition of $\bm V$  from the document suggests the following form of $\bm V$
 
\begin{equation}\label{V-1}
\bm V = (V_{ijpq}):=\begin{pmatrix}
V_{1111} & V_{1112} & V_{1113} &  0 &  0 & 0 \\
V_{1211} & V_{1212} & V_{1213} &  0 &  0 & 0 \\
V_{1311} & V_{1312} & V_{1313} &  0 &  0 & 0 \\
0 & 0 & 0 & V_{2121} &  V_{2122} & V_{2123} \\
0 &  0 & 0 & V_{2221} &  V_{2222} & V_{2223} \\
0 &  0 & 0 & V_{2321} &  V_{2322} & V_{2323}
\end{pmatrix}
\end{equation}
In such case, we are unsure about the positivity of $\bm V^{-1}$; such positivity is not guaranteed, which bothers us a bit. However, $\bm H$ can still be positive, but this will not be easy to establish, or, at least, it is not obvious to us.

\item On the other hand, if $\bm V$ is non-zero only when $i=p$ and $j=q$, then the positivity of $\bm{V}^{-1}$ can be guaranteed, and $\bm{V}$ is block diagonal.
\begin{equation}\label{V-2}
\bm V =  (V_{ijpq}) := \begin{pmatrix}
V_{1111} & 0 & 0 &  0 &  0 & 0 \\
0 & V_{1212} & 0 &  0 &  0 & 0 \\
0 & 0 & V_{1313} &  0 &  0 & 0 \\
0 & 0 & 0 & V_{2121} &  0 & 0 \\
0 &  0 & 0 & 0 &  V_{2222} & 0 \\
0 &  0 & 0 & 0 &  0 & V_{2323}
\end{pmatrix}
\end{equation}
\item {QUESTION: Which of the two definitions of $\bm V$ is the correct one (equation~\eqref{V-1} or equation~\eqref{V-2})}. 
\item {\bf ANSWER:} The $\bm V$ given in~\eqref{V-2} is the correct one. 

\item QUESTION(s) 1: What is the number of columns in $\bm{F}_{ijpq}$? Is it $8$ because of the two zero columns added? The $m$ in $\eta_{qm}$ is not the same as the number of blocks $m=2$ defined in the first item, right? Anyway, how is the number of rows and columns in $F_{ijpq}$ determined as a function of the sites and the age groups  ($m$ and $n$  from the first item above).

\item {\bf Answer:} $F_{ijpq}$ are the derivatives of the RHS functions of the model equations on page 1 with respect to the 8 variables in the vector \[
v_{ij}=( E_{ijs}, E_{ijp}, I_{ija}, I_{ijps}, I_{ijs}, I_{ijh}, I_{ijpm}, I_{ijm} ) 
\]
(see the top of page 2) evaluated at $v_{ij}=(0,0,\cdots, 0)$
and $S_{ij}=S^*_{ij}, V_{ij}=V^*_{ij}, P_{ij}=P^*_{ij}$.  The letter $m$ in subscripts of variables (e.g., $E$ or $I$) and parameters (e.g., $\eta$) are independent of the number of age groups ($n$) or sites ($m$).

\item Question 2: Does $\bm{V}$ depend on $\bm{\chi}$? It will be great to know either way. It will be easier to deal with the optimization problem (minimizing the entries of $\bm F$ if $\bm V$ is independent of $\bm \chi$, but this is not a show-stopper.

\item {\bf Answer:} The matrix $V$ does NOT depend on $\chi$.
 
\item Question 3: Which entries of $\bm F$ depend on $\bm \chi$. All? Some? How?

\item {\bf Answer:} Because all elements of $F$ depend on either $Y$ or $Z$, where $Y$ and $Z$ depend on $S^*/N=s^*$, or $V^*/N=v^*$, or $P^*/N=p^*$, which depend on $\chi$ as shown in line 5 on page 2, all non-zero elements of F depend on $\chi$.

\item {\bf Additional explanations:} In the attached pdf, I highlighted the items related to the v vector, $s^*$, $v^*$ etc. $Y_{ij}$ and $Z_{ij}$.
 
The constraint on $\chi$ is illustrated on page 4 (Equation (215)). But for the real model, the constraints consist of $m$ equations ($m$ is the number of sites) of the form:
\[
\chi_{i1}N_{i1} +\chi_{i2}N_{i2}+\ldots +\chi_{in}N_{in}=c_i,     \quad i=1:m.
\] 
Here, $N_{ij}$ is the total number of individuals of location  $i$ and age $j$. When you need these $N_{ij}$ values as well as other parameter values, we can ask John.
 
Although $Y$ and $Z$ do not depend on $(pq)$, the factor $c_{ijpq}$ in $F_{ij}$ depend of parameters associated with $p$ or $q$. Thees parameters $c_{ijpq}$ link populations of different ages and different sites.
 
\end{itemize}

\section{Code}

%\lstset{style=pystyle}

%\begin{lstlisting}
%\lstinputlisting[firstline=300,lastline=500]{Code/data_math_bio.py}
\lstinputlisting[language=Python,
       basicstyle=\ttfamily,
       keywordstyle=\color{blue},
       stringstyle=\color{red},
       commentstyle=\color{green}]{Code/z_math_bio.py}
%\end{lstlisting}
\end{document}
